\begin{enumerate}
    \item Benägenheten för lokförare att följa hastighetsrekommendationer från en C-DAS-klient riskerar att bli låg om tillämpningen av C-DAS-klienten görs oaktsamt.
    \item Utvecklare av en C-DAS-lösning bör beakta lokförarens perspektiv och tillämpa en applikation som inte distraherar föraren i onödan och som upprätthåller ett högt förtroendekapital hos lokförarkåren.
    \item Hastighetsrekommendationer bör återges först när tågfärden efter start har nått den övre hastighetsgränsen för sin journey profile och med siffror ange inom vilket spann/hastighetsfönster tåget kan röra sig (snarare än en exakt hastighet) eller om coasting (frirull) ska tillämpas.
    \item Hastighetsrekommendationer bör återges först när tågfärden efter start har nått den övre hastighetsgränsen för sin journey profile och med siffror ange inom vilket spann/hastighetsfönster tåget kan röra sig (snarare än en exakt hastighet) eller om coasting (frirull) ska tillämpas.
    \item För förare av godståg kan det vara svårare att uppfylla tågfärdens journey profile och behovet av frihet och manöverutrymme gällande tågets framförande är således större.
\end{enumerate}